We present a comprehensive validation study of the \emph{Gaussian encoder} used in
GaussianDiffusion, a diffusion model that operates in the latent space of 2-D Gaussian
splatting representations.  The encoder fits $K$ two-dimensional Gaussians to a greyscale
image by gradient descent, producing a compact $K\!\times\!7$ parameter tensor that
a downstream DDPM transformer subsequently learns to generate.

During development a critical \textbf{coordinate-inversion bug} was discovered: the
renderer (\texttt{F.affine\_grid}) inverts the translation coordinates, causing
Gaussian seeds placed at pixel positions $(x, y)$ to render on the opposite side of
the image.  Prior to the fix the encoder achieved only $26\,\mathrm{dB}$ mean PSNR
with $60\%$ dead Gaussians even after $3\,000$ epochs.  Three one-line changes to
the initialisation, dead-mask, and recycling routines corrected the convention, raising
mean PSNR to $\mathbf{44.2 \pm 3.7\,\mathrm{dB}}$ with $\mathbf{100\%}$ alive
Gaussians and convergence within ${\sim}790$ epochs on average.

We validate the encoder on MNIST across all ten digit classes, report per-digit PSNR
statistics, convergence dynamics, parameter distributions, spatial structure of the
learned Gaussians, K-sensitivity, and a loss-function ablation (MSE vs.\ L1\,+\,SSIM).
We further assess \emph{downstream readiness}: whether the encoded representations
are suitable as training data for a diffusion model, examining normalization coverage,
latent-space structure (PCA), encoding consistency, and the vestigial $\alpha$ parameter.
Our analysis confirms that the corrected encoder produces clean, normalizable, structurally
coherent representations that are ready for diffusion model training.
